% Converted from the authoritative recovery source, 2026-11-14.
\section{Recovering the influence proof for every exchangeable
law}\label{recovering-the-influence-proof-for-every-exchangeable-law}

Research deduction, 14 November 2026. This proves a direct generalized
influence inequality and then runs the original influence-potential
iteration. It applies to \textbf{all globally bounded polynomials}, not
only quantum acceptance functions. Uniform AA is the only non-elementary
input; the influence inequality is not inferred from a simulation
theorem.

\subsection{The theorem}\label{influence-recovery:the-theorem}

Let \(\mu\) be any exchangeable law on \(\{-1,1\}^n\), and let
\(f:\{-1,1\}^n\to[-1,1]\) have multilinear degree at most \(d\ge1\). Use
\(D_i f=(f(x^{i\leftarrow+1})-f(x^{i\leftarrow-1}))/2\) and
\(I_i^\mu(f)=\mathbb E_\mu(D_i f)^2\). Set \[
 v=\operatorname{Var}_\mu f,\qquad A_\mu(f)=\sum_i I_i^\mu(f),\qquad
 \Psi_\mu(f)=\operatorname{Var}_\mu f+A_\mu(f).
\] For a deterministic query tree \(\mathcal T\) with final transcript
\(H\), let \(\mu_H,f_H\) be its conditional law and restricted
polynomial on the remaining coordinates. Define its
\textbf{reveal-and-derivative influence} by \[
 \mathcal J_\mu(f;\mathcal T)
 =\Psi_\mu(f)-\mathbb E_H\Psi_{\mu_H}(f_H).                 \tag{1}
\] This equals the variance explained by the transcript plus the
expected cumulative derivative influences of the queried coordinates, as
proved below. Like the original derivative influence, it depends on the
polynomial's full-cube extension.

\textbf{Theorem.} There are absolute constants \(c,C>0\) such that
whenever \(v>0\), a deterministic tree of depth at most \[
 K(d,v)=\lceil C d^{18}v^{-7}\rceil
\] satisfies \[
 \boxed{\mathcal J_\mu(f;\mathcal T)\ge c v^{7/2}/d^7.}    \tag{2}
\] For large \(n\), the tree consists of a fixed prefix of queries
followed by one coordinate chosen from those answers. For small \(n\),
it queries the entire input. The result is uniform after arbitrary
positive-probability coordinate pinnings, since the remaining law stays
exchangeable.

The influence-potential iteration consequently gives a deterministic
\(O(d^{27}\epsilon^{-21})\)-query approximation with mean absolute error
\(\epsilon\). All constants are independent of dimension, temperature,
slice weights, full support, and minimum atom probabilities. In
particular, this recovers an influence proof for Curie--Weiss at every
temperature, including growing inverse temperatures and the
zero-temperature phase mixture. The exponents are not claimed optimal.
